\documentclass[aspectratio=169,11pt]{beamer}
\usetheme{Madrid}
\usecolortheme{default}
\setbeamertemplate{navigation symbols}{}
\setbeamertemplate{footline}[frame number]

\usepackage[utf8]{inputenc}
\usepackage[T1]{fontenc}
\usepackage[provide=*,italian]{babel}
\usepackage{amsmath,amssymb,mathtools}
\usepackage{bm}
\usepackage{array}

\usepackage{tikz}
\usetikzlibrary{calc,positioning,arrows.meta}

\newcommand{\tr}{\operatorname{tr}}
\newcommand{\C}{\mathbb{C}}
\newcommand{\R}{\mathbb{R}}
\newcommand{\D}{\mathbb{D}}
\newcommand{\M}{\mathbb{M}^{2,1}}
\newcommand{\HH}{\mathbb{H}}
\newcommand{\OD}{\mathcal{O}(\D)}
\newcommand{\EE}{\mathbf{E}}
\newcommand{\Amap}{\mathcal{A}}
\newcommand{\I}{\mathbf{I}}
\newcommand{\II}{\mathbf{II}}
\newcommand{\Iff}{\mathbf{I}}
\newcommand{\IIff}{\mathbf{II}}
\newcommand{\norm}[1]{\left\lVert #1 \right\rVert}

\title{Classificazione delle superfici space-like iperboliche\\ a curvatura media costante in spazio di Minkowski}
\subtitle{Tesi di Laurea in Matematica}
\author{\centering \textbf{Candidato:} Elia Andrea Perli\\
        \textbf{Relatore:} Prof.\ Francesco Bonsante}
\institute{Università degli Studi di Pavia}
\date{2026 -- 04 -- 28}

\begin{document}

\begin{frame}
  \titlepage
\end{frame}


% =================== SETTING GEOMETRICO ===================
\section{Setup geometrico}

\begin{frame}[t]{Lo spazio di Minkowski e le superfici space-like}
  \emph{Spazio di Minkowski}: $\M=\R^3$ con pseudometrica
  \[
    ds^2 \;=\; (dx^1)^2+(dx^2)^2-(dx^3)^2.
  \]
  Una superficie immersa $\iota:S\hookrightarrow\M$ è \alert{di tipo
  spazio} se la metrica indotta $\Iff:=\iota^* g_\M$ è Riemanniana.

  \vspace{2mm}
  \begin{block}{Forme fondamentali e curvature}
    Come nel caso euclideo, $S$ è equipaggiata con la \emph{prima
    forma fondamentale} $\Iff$ (metrica indotta) e la \emph{seconda
    forma fondamentale} $\IIff$ (variazione della normale unitaria),
    e si parla di \emph{curvatura media} $H=\tfrac12\tr_{\Iff}(\IIff)$ e
    \emph{curvatura di Gauss} $K$. Le forme $(\Iff,\IIff)$ soddisfano
    le equazioni di \alert{Gauss-Codazzi}.
  \end{block}

  \vspace{2mm}
  \emph{Esempio canonico:} grafico di $f:\R^2\to\R$. La condizione
  space-like si traduce in $\;|\nabla f|<1$, cioè $f$
  \alert{$1$-Lipschitz}: il piano tangente non è ``troppo verticale''.
\end{frame}


% =================== CMC E COMPLETEZZA ===================
\begin{frame}[t]{Superfici CMC complete}
  Ci concentriamo su superfici \alert{a curvatura media costante}
  (CMC) e \alert{complete}, dette \emph{tagli CMC}.

  \vspace{2mm}
  \begin{block}{Teorema (Cheng--Yau)}
    Per una superficie space-like CMC, sono equivalenti:
    \begin{itemize}\itemsep=0pt
      \item $S$ è un grafico intero su $\R^2$;
      \item la prima forma fondamentale $\Iff$ è una metrica
            Riemanniana \alert{completa}.
    \end{itemize}
  \end{block}

  \vspace{2mm}
  \emph{Analogo del teorema di Bernstein.}\;
  Una CMC con $K\equiv 0$ è un iperpiano $\Rightarrow$ non interessante;
  ci occuperemo del caso $K<0$ ovunque, detto \alert{iperbolico}.
  Per omotetia, ogni CMC con $H=\kappa>0$ corrisponde ad una con $H=1$.

  \vspace{3mm}
  \begin{alertblock}{Obiettivo della tesi}
    \centering
    Classificare le \emph{CMC-$1$ iperboliche} a meno di isometria di $\M$.
  \end{alertblock}
\end{frame}


% =================== COORDINATE ISOTERME ===================
\section{Riduzione del problema}

\begin{frame}[t]{Coordinate isoterme e parametrizzazione conforme}
  \begin{block}{Coordinate isoterme}
    Su una superficie Riemanniana orientata esistono coordinate locali
    $(x,y)$ in cui la metrica si scrive
    \[
      \Iff \;=\; e^{2u}\,(dx^2+dy^2),
      \qquad u:U\to\R.
    \]
    Posto $z=x+iy$, definiscono una struttura complessa.
  \end{block}

  \vspace{1mm}
  Una CMC-$1$ iperbolica $S$ è semplicemente connessa (grafico intero):
  per il \alert{teorema di uniformizzazione} esiste un diffeomorfismo
  conforme $\sigma:\D\to S$.

  \vspace{1mm}
  In tali coordinate, su tutto $\D$:
  \[
    \Iff \;=\; e^{2u}|dz|^2,
    \qquad
    \IIff \;=\; h_{ij}\,dx^i\,dx^j.
  \]
\end{frame}


% =================== INCOGNITE DEL PROBLEMA ===================
\begin{frame}[t]{Le incognite del problema}
  Le \alert{quattro funzioni incognite} $u,h_{11},h_{12},h_{22}$ sono
  legate dalle relazioni geometriche.

  \vspace{2mm}
  \begin{columns}[T,onlytextwidth]
    \column{0.48\textwidth}
    \begin{block}{Vincolo CMC-$1$}
      \centering
      Da $H=\tfrac12\tr_{\Iff}(\IIff)=1$:
      \[
        h_{11}+h_{22} \;=\; 2e^{2u}.
      \]
      $h_{22}$ è dunque determinato.
    \end{block}
    \column{0.48\textwidth}
    \begin{block}{Incognite indipendenti}
      \centering
      Restano \alert{tre} funzioni:
      \[
        u,\quad h_{11},\quad h_{12}\,:\,\D\to\R.
      \]
    \end{block}
  \end{columns}

  \vspace{3mm}
  Resta da imporre le equazioni di \alert{Gauss-Codazzi}, che
  garantiscono che $(\Iff,\IIff)$ provengano effettivamente da una
  superficie immersa in $\M$ (\emph{teorema di Bonnet}).
\end{frame}


% =================== CODAZZI E FUNZIONE OLOMORFA ===================
\begin{frame}[t]{Codazzi: una funzione olomorfa sul disco}
  Definiamo la funzione $f:\D\to\C$ ponendo
  \[
    f(z) \;:=\; \tfrac{h_{11}-h_{22}}{2}\;-\;i\,h_{12}
        \;=\; (h_{11}-e^{2u})\;-\;i\,h_{12}.
  \]

  \vspace{1mm}
  \begin{columns}[T,onlytextwidth]
    \column{0.10\textwidth}\hfill
    \column{0.80\textwidth}
    \begin{block}{Proposizione (traduzione di Codazzi)}
      \centering
      Le equazioni di Codazzi per $\IIff$ rispetto ad $\Iff=e^{2u}|dz|^2$ sono
      equivalenti a
      \[
        f \text{ è \alert{olomorfa} su } \D.
      \]
    \end{block}
    \column{0.10\textwidth}\hfill
  \end{columns}

  \vspace{2mm}
  \emph{Idea.} In coordinate isoterme, dopo cancellazione dei
  simboli di Christoffel di $\Iff$, le equazioni di Codazzi sulla
  parte traceless di $\IIff$ si riducono alle equazioni di
  \alert{Cauchy-Riemann} per $f$.

  \vspace{2mm}
  \emph{Conseguenza.} I dati $(h_{11},h_{12})$ si recuperano dalla
  parte reale ed immaginaria di $f$. Resta da determinare il fattore
  conforme $u$.
\end{frame}


% =================== GAUSS COME PDE ===================
\begin{frame}[t]{Gauss: una PDE non lineare per il fattore conforme}
  Dato $f$ olomorfa, l'equazione di Gauss diventa un'equazione
  \alert{solo per $u$}.

  \vspace{1mm}
  Per ottenerne una forma analiticamente trattabile, si scrive $\Iff$
  rispetto alla \emph{metrica di Poincaré}
  $ds_p^2 = \frac{4}{(1-|z|^2)^2}|dz|^2$ (completa, di curvatura $-1$):
  \[
    \Iff \;=\; e^{2w}\,ds_p^2,
    \qquad
    w:\D\to\R.
  \]

  \vspace{1mm}
  \begin{columns}[T,onlytextwidth]
    \column{0.10\textwidth}\hfill
    \column{0.80\textwidth}
    \begin{block}{Proposizione (traduzione di Gauss)}
      \centering
      L'equazione di Gauss è equivalente alla PDE
      \[
        \boxed{\;\Delta_p w \;=\; e^{2w}\;-\;\norm{f}^2 e^{-2w}\;-\;1\;}
        \quad\text{su }\D,
      \]
      dove $\Delta_p$ è il laplaciano di Poincaré e $\norm{f}$ la
      norma di $f$ rispetto a $ds_p^2$.
    \end{block}
    \column{0.10\textwidth}\hfill
  \end{columns}
\end{frame}


% =================== STRATEGIA DI INVERSIONE ===================
\begin{frame}[t]{Strategia di inversione e enunciato}
  \emph{Schema.} Data una funzione olomorfa $f:\D\to\C$:
  \begin{enumerate}\itemsep=0pt
    \item si cerca $w$ soluzione della PDE, con $e^{2w}ds_p^2$ completa;
    \item si definiscono $h_{11},h_{12}$ tramite $f$ (Codazzi gratis);
    \item dalla CMC-$1$, $h_{22}=2e^{2w}\rho_p^{-2}-h_{11}$, dove
          $\rho_p^{-2}=\frac{(1-|z|^2)^2}{4}$;
    \item per il \emph{teorema di Bonnet}, esiste $S\subset\M$ con
          tali $\Iff,\IIff$, unica a meno di isometria.
  \end{enumerate}

  \vspace{2mm}
  Definendo
  $\EE = \{\text{CMC-$1$ iperboliche}\}/\text{isometrie di }\M$ e
  $\OD = \{f:\D\to\C \text{ olomorfa}\}/\text{mappe conformi di }\D$,
  la costruzione produce
  $\Amap:\EE\to\OD,\;[S]\mapsto[f]$.

  \vspace{2mm}
  \begin{columns}[T,onlytextwidth]
    \column{0.10\textwidth}\hfill
    \column{0.80\textwidth}
    \begin{block}{Teorema di classificazione (Wan, 1992)}
      \centering
      L'applicazione $\Amap:\EE\to\OD$ è
      \alert{ben definita e biunivoca}.
    \end{block}
    \column{0.10\textwidth}\hfill
  \end{columns}
\end{frame}


% =================== IL PROBLEMA ANALITICO ===================
\begin{frame}[t]{Il problema analitico}
  La biunivocità di $\Amap$ si riduce a:

  \vspace{2mm}
  \begin{columns}[T,onlytextwidth]
    \column{0.07\textwidth}\hfill
    \column{0.86\textwidth}
    \begin{block}{Problema}
      Data $f:\D\to\C$ olomorfa, esiste \alert{un'unica}
      $w\in C^\infty(\D)$ tale che
      \[
        \Delta_p w \;=\; e^{2w}\;-\;\norm{f}^2 e^{-2w}\;-\;1
        \quad\text{su }\D,
      \]
      con $e^{2w}ds_p^2$ \alert{metrica completa} su $\D$.
    \end{block}
    \column{0.07\textwidth}\hfill
  \end{columns}

  \vspace{3mm}
  \begin{itemize}\itemsep=2pt
    \item Non è un problema di Dirichlet: \alert{nessun dato al bordo}.
    \item La condizione di completezza richiede a $w$ di
          \emph{divergere opportunamente} all'avvicinarsi a $\partial\D$.
    \item Affronteremo nell'ordine: prima l'\alert{esistenza}, poi
          l'\alert{unicità}.
  \end{itemize}
\end{frame}


% =================== ESISTENZA - METODO ===================
\section{Esistenza}

\begin{frame}[t]{Esistenza: il metodo delle sotto/sovra-soluzioni}
  La tecnica è classica: trovare due funzioni $\underline\psi\le\overline\psi$,
  rispettivamente \alert{sotto-soluzione} e \alert{sovra-soluzione}
  della PDE, e dedurne una soluzione $w$ con $\underline\psi\le w\le\overline\psi$.

  \vspace{2mm}
  \begin{block}{Teorema (metodo monotono)}
    Su una varietà Riemanniana completa, se $F(x,u)\in C^\infty(M\times\R)$
    con $\partial_u F>0$ ed esistono
    $\underline\psi\le\overline\psi\in C^0(M)\cap H^1(M)$ con
    $\Delta\underline\psi\ge F(x,\underline\psi)$,
    $\Delta\overline\psi\le F(x,\overline\psi)$,
    allora esiste $u\in C^\infty(M)$ con
    $\Delta u=F(x,u)$ e $\underline\psi\le u\le\overline\psi$.
  \end{block}

  \vspace{2mm}
  Nel nostro caso $F(x,w)=e^{2w}-\norm{f(x)}^2 e^{-2w}-1$, con
  $\partial_w F = 2(e^{2w}+\norm{f}^2 e^{-2w}) > 0$.
\end{frame}


% =================== ESISTENZA - CASO LIMITATO ===================
\begin{frame}[t]{Esistenza: il caso $\norm{f}$ limitata}
  \emph{Ipotesi:} $\norm{f}$ \alert{limitata}, $a:=\sup_\D\norm{f}<\infty$.
  Cerchiamo sotto/sovra-soluzioni costanti.

  \vspace{2mm}
  \begin{columns}[T,onlytextwidth]
    \column{0.48\textwidth}
    \begin{block}{Sotto-soluzione: $\underline\psi\equiv 0$}
      \centering
      \[
        \Delta_p 0 - F(x,0) \;=\; \norm{f}^2 \;\ge\; 0.\qquad\checkmark
      \]
    \end{block}
    \column{0.48\textwidth}
    \begin{block}{Sovra-soluzione: $\overline\psi\equiv C$}
      \centering
      Posto $X=e^{2C}$, basta $X^2-X-a^2\ge 0$:
      \[
        C \;=\; \tfrac12\log\!\left(\tfrac{1+\sqrt{1+4a^2}}{2}\right).
      \]
    \end{block}
  \end{columns}

  \vspace{3mm}
  Per il metodo monotono otteniamo $w\in C^\infty(\D)$ con
  \[
    0 \;\le\; w \;\le\; C
    \quad\Longrightarrow\quad
    1 \;\le\; e^{2w} \;\le\; \tfrac{1+\sqrt{1+4a^2}}{2}.
  \]

  \vspace{1mm}
  \begin{alertblock}{Completezza, gratis}
    \centering
    $w \ge 0\;\Rightarrow\; e^{2w}ds_p^2 \ge ds_p^2$.\\
    Poiché $ds_p^2$ è completa su $\D$, anche $e^{2w}ds_p^2$ lo è.
  \end{alertblock}
\end{frame}


% =================== UNICITÀ - PRINCIPIO DEL MASSIMO ===================
\section{Unicità}

\begin{frame}[t]{Unicità: il principio del massimo generalizzato}
  Su domini non compatti la mancanza di un punto di massimo si rimpiazza
  con \emph{quasi-massimi}, sfruttando la completezza della metrica.

  \vspace{2mm}
  \begin{columns}[T,onlytextwidth]
    \column{0.07\textwidth}\hfill
    \column{0.86\textwidth}
    \begin{block}{Principio del massimo generalizzato (Omori-Yau)}
      Sia $(M,g)$ Riemanniana completa con curvatura di Ricci limitata
      dal basso, $\eta\in C^2(M)$ limitata superiormente. Allora
      esiste $(p_n)\subset M$ con
      \[
        \eta(p_n)\to\sup_M\eta,
        \qquad
        |\nabla_g\eta|(p_n)\to 0,
        \qquad
        \limsup_n \Delta_g\eta(p_n)\le 0.
      \]
    \end{block}
    \column{0.07\textwidth}\hfill
  \end{columns}

  \vspace{2mm}
  \begin{itemize}\itemsep=2pt
    \item Nel caso generale la dimostrazione è \emph{tecnica} (stime
          del laplaciano della distanza tramite Hessiana).
    \item Nel nostro caso (\alert{dimensione 2}, curvatura non
          positiva) bastano argomenti più semplici basati sui
          confronti di Hessiana.
  \end{itemize}
\end{frame}


\begin{frame}[t]{Unicità: applicazione alla PDE}
  Siano $w_1,w_2$ soluzioni del problema. Poniamo $\eta = w_2-w_1$ e
  dimostriamo $\eta\le 0$ (la disuguaglianza inversa segue per simmetria).

  \vspace{2mm}
  \begin{columns}[T,onlytextwidth]
    \column{0.07\textwidth}\hfill
    \column{0.86\textwidth}
    \begin{block}{Identità chiave}
      Posto $|\mu|^2=\norm{f}^2 e^{-4 w_1}\le 1$
      (la condizione di positività),
      \[
        \Delta_{g_1}\eta
        \;=\; e^{2\eta}-|\mu|^2 e^{-2\eta}-1+|\mu|^2
        \;\ge\; e^{2\eta}-e^{-2\eta}-1
        \;=:\; F(\eta).
      \]
      $F$ è strettamente crescente con $F(0)=0$.
    \end{block}
    \column{0.07\textwidth}\hfill
  \end{columns}

  \vspace{2mm}
  Applicando il principio del massimo a $\eta$ sulla varietà completa
  $(\D,g_1=e^{2w_1}ds_p^2)$:
  \[
    \exists\,(p_n):\;
    \eta(p_n)\to\bar\eta:=\sup\eta,\quad
    \limsup_n\Delta_{g_1}\eta(p_n)\le 0
  \]
  e passando al limite, $F(\bar\eta)\le 0$, cioè $\bar\eta\le 0$.
  Per simmetria, $w_1=w_2$.\quad\alert{Unicità.}
\end{frame}


% =================== MOTIVAZIONE / CONCLUSIONI ===================
\section{Conclusione}

\begin{frame}[t]{Una motivazione: mappe armoniche $\D\to\D$}
  Una motivazione importante per lo studio delle CMC-$1$ in $\M$ viene
  dalla \alert{teoria delle mappe armoniche}.

  \vspace{2mm}
  \begin{itemize}\itemsep=4pt
    \item Per ogni CMC-$1$ iperbolica $S\subset\M$, la sua mappa di
          Gauss letta in coordinate conformi è una \alert{mappa armonica}
          $u:\D\to\D$.
    \item Viceversa, ogni mappa armonica $\D\to\D$ proviene da una
          tale CMC-$1$.
    \item La classificazione $\Amap$ fornisce dunque un dizionario tra
          \emph{geometria delle CMC} in $\M$ e
          \emph{teoria delle mappe armoniche} sul disco;
          è il punto di contatto con lo \emph{spazio di Teichmüller
          universale} (Wan, 1992).
  \end{itemize}
\end{frame}


\begin{frame}[t]{Conclusioni}
  \begin{columns}[T,onlytextwidth]
    \column{0.07\textwidth}\hfill
    \column{0.86\textwidth}
    \begin{block}{Teorema di classificazione}
      \centering
      $\Amap:\EE\to\OD,\;[S]\mapsto[f]\;$
      è \alert{ben definita e biunivoca}.
    \end{block}
    \column{0.07\textwidth}\hfill
  \end{columns}

  \vspace{3mm}
  \textbf{In sintesi.}
  \begin{itemize}\itemsep=2pt
    \item In coordinate isoterme, l'equazione di Codazzi diventa
          l'\emph{olomorfia} di una funzione $f:\D\to\C$;
    \item l'equazione di Gauss diventa la PDE non lineare
          $\Delta_p w = e^{2w}-\norm{f}^2 e^{-2w}-1$;
    \item la biunivocità di $\Amap$ corrisponde a esistenza ed unicità
          di $w$, con $e^{2w}ds_p^2$ \emph{completa}.
  \end{itemize}

  \vspace{2mm}
  \textbf{Strumenti analitici.}
  \begin{itemize}\itemsep=2pt
    \item Esistenza: \alert{metodo delle sotto/sovra-soluzioni}.
    \item Unicità: \alert{principio del massimo generalizzato}.
  \end{itemize}
\end{frame}


\begin{frame}
  \vfill
  \begin{center}
    \Huge Grazie per l'attenzione.
  \end{center}
  \vfill
\end{frame}


\begin{frame}[allowframebreaks]{Riferimenti essenziali}
\footnotesize
\begin{thebibliography}{9}
\bibitem{Wan} T.~Y.~H.~Wan, \emph{Constant mean curvature surface,
  harmonic maps, and universal Teichmüller space}, J.\ Differential
  Geom.\ 35 (1992), 643--657.
\bibitem{Jost} J.~Jost, \emph{Compact Riemann Surfaces}, Springer, 2006.
\bibitem{doCarmo} M.~P.\ do Carmo, \emph{Differential Geometry of
  Curves and Surfaces}, Dover, 2016.
\end{thebibliography}
\end{frame}

\end{document}
